%----------------------------------------------------------------------------------------
% Research design
%----------------------------------------------------------------------------------------

\section{Use Cases} 
\subsection{Trading on Cryptocurrency Exchanges}
USB can be used for trading on various cryptocurrency exchanges, where it can offer advantages over traditional fiat-backed stablecoins. 
With its peg to the USD value and its use of Bitcoin inverse perpetual swaps for hedging, USB offers a stable and reliable trading pair for investors looking to hedge against market volatility. 
Its decentralized and censorship-resistant nature also makes it attractive to those who value privacy and security in their trading activities.\newline
USB can also be traded on decentralized exchanges (DEXs), where its use of Bitcoin inverse perpetual swaps can help to provide liquidity and maintain its peg to the USD value. 
Its decentralized and censorship-resistant nature also makes it a more attractive option for those who prefer to avoid centralized exchanges or who are located in countries with stricter regulations on cryptocurrency trading.

\subsection{Remittance Payments}
USB can be used as a reliable and cost-effective solution for remittance payments.
Traditional methods of sending money across borders are often expensive and can take several days to process. 
With USB, remittances can be sent quickly and at a low cost, as there are no intermediaries involved in the process. 
Additionally, the stable value of USB means that the recipient can be sure they will receive the full intended amount without worrying about fluctuations in exchange rates. 
This can greatly benefit individuals and businesses who rely on cross-border transactions for their livelihoods.
\subsection{E-commerce Transactions}
USB can also be used as a payment method for e-commerce transactions.
E-commerce platforms and merchants can easily integrate USB payments into their checkout process, allowing for fast and secure transactions.
The stable value of USB provides a reliable payment method for merchants, reducing the volatility risks associated with traditional cryptocurrency payment methods. 
Additionally, the low transaction fees associated with USB payments make it an attractive payment method for merchants looking to reduce payment processing costs.
With the added benefits of decentralization, censorship-resistance, and privacy, USB provides an ideal payment solution for e-commerce transactions.

\subsection{Store of Value for Long-term Holding}
USB is designed not only for day-to-day transactions but also as a store of value for long-term holding.
Traditional fiat currencies can be subject to inflation and devaluation, leading to decreased purchasing power over time. USB, on the other hand, is pegged to the value of the US dollar, which has historically been a relatively stable currency.\newline
By using inverse perpetual swaps for hedging, USB can maintain its peg to the US dollar without being subject to the same inflationary pressures as fiat currencies.
This makes it a potentially attractive option for those seeking to hold assets for the long term and preserve their value.\newline
Furthermore, USB's low counterparty risk and transparency regarding asset backing can provide added peace of mind for those looking for a secure and reliable store of value. 
With its decentralized nature, USB can also offer protection against potential government or regulatory interference in traditional fiat-based systems.
The additional yield generated from USB's hedging mechanisms (60\% for ETHUSD and 16\% for BTCUSD in 2024) further enhances its appeal as a store of value, combining stability with growth potential. It is important to note that past performance is not indicative of future results, and these figures should not be considered a guarantee of future returns. 

